\documentclass[11pt]{article}
\usepackage[margin=1in]{geometry}
\usepackage{amsmath,amssymb,amsthm}
\usepackage{graphicx}
\usepackage[hidelinks]{hyperref}
\setlength{\emergencystretch}{2em}

\newtheorem{theorem}{Theorem}
\newtheorem{lemma}{Lemma}
\newtheorem{corollary}{Corollary}
\newtheorem{remark}{Remark}

\newcommand{\C}{\mathbb C}
\newcommand{\Bcal}{\mathcal B}
\newcommand{\Pcal}{\mathcal P}

\title{Pure-State Pi-Bases Obstruct the Negative Von Neumann Cases}
\author{Run \texttt{fa\_banach\_001}, lane 19}
\date{June 17, 2026}

\begin{document}
\maketitle

\section*{Source Problem}

The source paper is Jinghao Huang, Karimbergen Kudaybergenov, and Rui Liu,
\emph{The ball-covering property of von Neumann algebras and noncommutative
symmetric spaces}, arXiv:2311.14257.  It proves that a von Neumann algebra
has BCP, equivalently UBCP, if and only if it is atomic and separably
represented.  The same paper asks for a broader noncommutative
\(C^*\)-algebra characterization:

\begin{center}
\includegraphics[width=0.92\textwidth]{figures/open_problem_crop.png}
\end{center}

\section*{Statement}

For a \(C^*\)-algebra \(A\), let \(\Pcal(A)\) denote its pure state space,
with the relative weak-star topology inherited from \(A^*\).

\begin{theorem}
\label{thm:main}
Let \(M\) be a von Neumann algebra.  If \(\Pcal(M)\) has a countable
\(\pi\)-base, then \(M\) is atomic and separably represented.  Consequently
\(M\) has BCP and UBCP by the source paper's von Neumann classification.
\end{theorem}

\begin{corollary}
\label{cor:negative}
If a von Neumann algebra \(M\) fails BCP, then \(\Pcal(M)\) has no countable
\(\pi\)-base.
\end{corollary}

\section*{Idea of the Proof}

The proof is by contrapositive and splits exactly along the source
classification.

If \(M\) has an atomless direct summand, all pure states on that summand are
singular.  Given countably many proposed pure-state open sets, choose one
pure state from each.  A common-projection lemma for countably many singular
states gives a nonzero projection \(p\) annihilating all of them.  But the
set of pure states satisfying \(\psi(p)>1/2\) is a nonempty open set, and
none of the proposed open sets is contained in it.

If \(M\) is atomic but not separably represented, then either it has
uncountably many atomic summands, or one summand is \(\Bcal(H)\) for a
nonseparable Hilbert space \(H\).  In both cases one obtains uncountably many
pairwise disjoint nonempty open subsets of the pure state space, which rules
out a countable \(\pi\)-base.

\section*{Atomless Summands}

\begin{lemma}[common projection for singular pure states]
\label{lem:common}
Let \(N\) be a nonzero von Neumann algebra and let
\((\varphi_n)\) be countably many singular states on \(N\).  Then there is a
nonzero projection \(p\in N\) such that
\[
        \varphi_n(p)=0\qquad(n\ge1).
\]
\end{lemma}

\begin{proof}
Fix a normal state \(\omega\) on \(N\).  A singular positive functional
\(\varphi\) admits an increasing net of projections \(e_i\uparrow1\) such
that \(\varphi(e_i)=0\).  Since \(\omega\) is normal, for every
\(\varepsilon>0\) there is a projection \(e\) with
\[
        \varphi(e)=0,\qquad \omega(1-e)<\varepsilon.
\]
Apply this to \(\varphi_n\) with \(\varepsilon_n=2^{-n-2}\), obtaining
projections \(e_n\) with
\[
        \varphi_n(e_n)=0,\qquad \omega(1-e_n)<2^{-n-2}.
\]
Let
\[
        p=\bigwedge_{n=1}^{\infty} e_n.
\]
Then
\[
        \omega(1-p)
        =
        \omega\Big(\bigvee_n(1-e_n)\Big)
        \le
        \sum_n\omega(1-e_n)
        <
        \frac12.
\]
Thus \(p\ne0\).  Since \(p\le e_n\), positivity gives
\(\varphi_n(p)=0\) for every \(n\).
\end{proof}

\begin{lemma}
\label{lem:atomless}
If \(N\) is a nonzero atomless von Neumann algebra, then \(\Pcal(N)\) has no
countable \(\pi\)-base.
\end{lemma}

\begin{proof}
In an atomless von Neumann algebra there are no normal pure states, because
normal pure states are exactly vector states supported on minimal
projections.  Thus every pure state on \(N\) is singular.

Let \((U_n)\) be any countable family of nonempty open subsets of
\(\Pcal(N)\).  Choose \(\varphi_n\in U_n\).  By Lemma~\ref{lem:common}, there
is a nonzero projection \(p\in N\) with \(\varphi_n(p)=0\) for every \(n\).

Choose a pure state \(\psi_0\) on the nonzero unital \(C^*\)-algebra
\(pNp\).  Its canonical extension
\[
        \psi(x)=\psi_0(pxp)\qquad(x\in N)
\]
is a pure state of \(N\), since pure states of hereditary subalgebras extend
purely to the ambient algebra.  Moreover \(\psi(p)=1\).

Therefore
\[
        W=\{\psi\in\Pcal(N):\psi(p)>1/2\}
\]
is a nonempty open subset of \(\Pcal(N)\).  But \(\varphi_n\in U_n\setminus W\)
for every \(n\).  Hence no \(U_n\) is contained in \(W\), so \((U_n)\) cannot
be a \(\pi\)-base.  Since the countable family was arbitrary, \(\Pcal(N)\)
has no countable \(\pi\)-base.
\end{proof}

\section*{Atomic Nonseparable Summands}

\begin{lemma}
\label{lem:disjoint}
If a topological space has uncountably many pairwise disjoint nonempty open
sets, then it has no countable \(\pi\)-base.
\end{lemma}

\begin{proof}
If \((V_n)\) were a countable \(\pi\)-base and \((O_i)_{i\in I}\) were
pairwise disjoint nonempty open sets, then for each \(i\) some \(V_{n(i)}\)
would satisfy \(V_{n(i)}\subset O_i\).  Distinct \(i\)'s require distinct
\(n(i)\)'s, because the \(O_i\)'s are disjoint and \(V_{n(i)}\ne\emptyset\).
Thus \(I\) would be countable.
\end{proof}

\begin{lemma}
\label{lem:atomic}
Let \(M\) be an atomic von Neumann algebra which is not separably represented.
Then \(\Pcal(M)\) has no countable \(\pi\)-base.
\end{lemma}

\begin{proof}
Write
\[
        M\cong \prod_{i\in I}\Bcal(H_i).
\]
If \(I\) is uncountable, let \(z_i\) be the central projection onto the
\(i\)-th summand.  The sets
\[
        O_i=\{\varphi\in\Pcal(M):\varphi(z_i)>1/2\}
\]
are pairwise disjoint nonempty open subsets of \(\Pcal(M)\).  By
Lemma~\ref{lem:disjoint}, \(\Pcal(M)\) has no countable \(\pi\)-base.

If \(I\) is countable but \(M\) is not separably represented, then some
summand is \(\Bcal(H)\) with \(H\) nonseparable.  It is enough to show that
\(\Pcal(\Bcal(H))\) has no countable \(\pi\)-base, since this pure-state
space is a clopen subspace of \(\Pcal(M)\).

Choose an uncountable orthonormal family \((\xi_\alpha)_{\alpha\in\Gamma}\)
in \(H\), and let \(p_\alpha\) be the rank-one projection onto
\(\C\xi_\alpha\).  The sets
\[
        W_\alpha
        =
        \{\varphi\in\Pcal(\Bcal(H)):\varphi(p_\alpha)>1/2\}
\]
are nonempty because they contain the vector states \(\omega_{\xi_\alpha}\).
They are open.  They are pairwise disjoint: if \(\alpha\ne\beta\), then
\(p_\alpha+p_\beta\le1\), so no state can take value \(>1/2\) on both
\(p_\alpha\) and \(p_\beta\).  Lemma~\ref{lem:disjoint} applies again.
\end{proof}

\section*{Proof of the Main Result}

\begin{proof}[Proof of Theorem~\ref{thm:main}]
We prove the contrapositive.  Suppose \(M\) is not atomic and separably
represented.

If \(M\) has a nonzero atomless central summand \(N=zM\), then
\[
        \Pcal(N)=\{\varphi\in\Pcal(M):\varphi(z)=1\}
\]
is a clopen subspace of \(\Pcal(M)\).  By Lemma~\ref{lem:atomless},
\(\Pcal(N)\) has no countable \(\pi\)-base.  Therefore \(\Pcal(M)\) has no
countable \(\pi\)-base, because every open subspace of a space with a
countable \(\pi\)-base inherits one.

If \(M\) is atomic but not separably represented, then
Lemma~\ref{lem:atomic} gives the same conclusion.

Thus \(\Pcal(M)\) having a countable \(\pi\)-base forces \(M\) to be atomic
and separably represented.  The source paper's von Neumann theorem then gives
BCP and UBCP.
\end{proof}

\begin{proof}[Proof of Corollary~\ref{cor:negative}]
The source paper proves that a von Neumann algebra fails BCP exactly when it
is not atomic and separably represented.  Apply Theorem~\ref{thm:main}.
\end{proof}

\section*{Scope and Limitations}

This packet proves a one-way compatibility theorem for the pure-state
\(\pi\)-base route in the von Neumann class:
\[
        \Pcal(M)\text{ has a countable }\pi\text{-base}
        \quad\Longrightarrow\quad
        M\text{ has UBCP}.
\]
It does not prove the converse for all atomic separably represented von
Neumann algebras, and it does not settle arbitrary \(C^*\)-algebras.  Its
role is to show that the source's von Neumann negative cases are already
detected at the level of pure-state \(\pi\)-bases.

\section*{Novelty and Search Bounds}

Local indexes were searched for arXiv:2311.14257, von Neumann BCP, pure
state space, \(\pi\)-base, singular pure state, atomless von Neumann algebra,
and atomic nonseparable representation.  The source paper contains the von
Neumann BCP/UBCP classification, and previous lane-19 packets contain the
countable-capture refinement.  No prior local packet recorded this
pure-state \(\pi\)-base obstruction for the negative von Neumann cases.

\section*{Human Review Focus}

Reviewers should check:
\begin{enumerate}
    \item the common-projection lemma for countably many singular states;
    \item the assertion that atomless von Neumann algebras have no normal pure
          states;
    \item the pure extension from \(pNp\) to \(N\);
    \item the inheritance of countable \(\pi\)-bases by open subspaces;
    \item the uncountable disjoint-open-family arguments in the atomic
          nonseparable cases.
\end{enumerate}

\begin{thebibliography}{9}

\bibitem{HKL}
Jinghao Huang, Karimbergen Kudaybergenov, and Rui Liu,
\emph{The ball-covering property of von Neumann algebras and noncommutative
symmetric spaces},
arXiv:2311.14257, 2023.

\bibitem{Blackadar}
Bruce Blackadar,
\emph{Operator Algebras: Theory of \(C^*\)-Algebras and Von Neumann
Algebras},
Encyclopaedia of Mathematical Sciences 122, Springer, 2006.

\bibitem{lane19vna}
Run \texttt{fa\_banach\_001}, lane 19,
\emph{Countable Capture for Von Neumann Algebras},
solution packet
\texttt{2311.14257\_von\_neumann\_capture\_classification}, 2026.

\end{thebibliography}

\end{document}
